% A proved intrinsic conditional lemma; geometric realization is not asserted here.
\paragraph{A fixed-node correction for a two-component exceptional surface
(intrinsic auxiliary lemma).}
Let $k$ be algebraically closed of characteristic $p>0$ and let
$E=S\cup R$ be a projective SNC surface with two smooth connected
components and smooth connected double curve $C$. Suppose that
$S$ is a smooth projective surface, that $R$ is a blow-up of
$\mathbf P^2$ at distinct points of a smooth plane model of $C$,
and that $C\subset R$ is its strict transform. Write these points
as $z_1,\ldots,z_v$, where $v>0$, and the corresponding exceptional
curves as $e_1,\ldots,e_v$. Thus
\[
 \mathcal O_R(e_j)|_C=\mathcal O_C(z_j).
\]
Assume that there are an ample line bundle $A$ on $S$, an integer
$a>0$, a section $h\in H^0(S,A^a)$ cutting out $C$, and a section
$g\in H^0(S,A^b)$ such that
\[
                \operatorname{div}(g|_C)=z_1+\cdots+z_v
                                                        \tag{FN.1}
\]
is reduced. Let $S(E)=F_\gamma^2K_0(E)_{\mathbf Q}$ and define
$J_K(E)$ by actual smooth projective curve families as in EP.
Then component degrees induce an isomorphism
\begin{equation}
                    S(E)/J_K(E)\simeq\mathbf Q^2.      \tag{FN.2}
\end{equation}
Consequently, if this $E$ is the full exceptional surface of data
satisfying EP and has a positive exceptional Cartier combination
with ample negative restriction, then $B^2(X)_{\mathbf Q}=0$.
The last implication is conditional on that geometric realization.

\emph{1. Move on the first component with fixed boundary support.}
Fix $x,y\in(S\setminus C)(k)$. Choose $m$ so large that the
sections of $A^{mb-a}$ separate the prescribed values and first
jets at $x,y$, and supply all required first jets away from $C$.
In the affine linear system
\begin{equation}
             T=\{g^m+hq=0\},\qquad
               q\in H^0(S,A^{mb-a}),                 \tag{FN.3}
\end{equation}
impose the two linear conditions that $T$ contain $x$ and $y$.
Choose $q(z_j)\ne0$ for every $j$. A general such section defines
a smooth connected curve $T$. Here is the precise elementary
Bertini argument. Away from $C$, multiplication by the nonzero
section $h$ preserves first-jet generation, so the incidence of
a singular point has codimension three over a two-dimensional
surface. The same estimate after imposing the two values, and
the remaining free derivative choices at $x,y$, avoids all
singularities there. On $C$ the only zeroes are $z_j$. For $m>1$
the differential there is $q(z_j)\,dh$, which is nonzero, since
$C$ is smooth. For $m=1$ its restriction to the tangent of $C$
is $dg|_C\ne0$. Finally, a sufficiently ample smooth divisor
on a connected smooth projective surface is connected, for
example by Serre duality and Serre vanishing applied to
$0\to A^{-mb}\to\mathcal O_S\to\mathcal O_T\to0$.
The chosen curve satisfies the exact scheme-theoretic intersection
\begin{equation}
                     T\cap C=m\sum_jz_j.             \tag{FN.4}
\end{equation}
Denote its parameter points over $z_j$ by $t_j\in T(k)$.
The points $x,y$ are different from all the $t_j$.

The graph of the inclusion $T\hookrightarrow S$ gives a perfect
class on $S\times T$. Replace it by its rational Adams-weight-two
projection, denoted $\alpha_S$. Its restrictions at $x,y$ are
$[\mathcal O_x]$ and $[\mathcal O_y]$. On the smooth surface
$C\times T$, refined intersection and (FN.4) give
\begin{equation}
      \alpha_S|_{C\times T}
           =m\sum_j[\mathcal O_{z_j}]\boxtimes
                             [\mathcal O_{t_j}].      \tag{FN.5}
\end{equation}
The length-$m$ local intersection is a zero-dimensional sheaf
on the regular surface $C\times T$, so its perfect $K_0$ class
is exactly $m$ times the residue-field class.

On $R\times T$ take the weight-two projection of
\[
             m\sum_j[\mathcal O_{e_j}]\boxtimes
                                    [\mathcal O_{t_j}],
\]
and call it $\alpha_R$. Transversality of $e_j$ and $C$ shows
that its restriction is the right side of (FN.5), and its
restrictions at $x,y$ are zero.

Rational closed-union descent in positive characteristic, as
used in FF and TS, gives the exact Mayer--Vietoris sequence
\begin{equation}
 K_1(C\times T)_{\mathbf Q}\longrightarrow K_0(E\times T)_{\mathbf Q}
 \longrightarrow K_0(S\times T)_{\mathbf Q}\oplus
                   K_0(R\times T)_{\mathbf Q}
 \longrightarrow K_0(C\times T)_{\mathbf Q}.          \tag{FN.6}
\end{equation}
Thus $(\alpha_S,\alpha_R)$ has an actual lift. Apply the
rational weight-two projector on the projective threefold
$E\times T$ to that lift. The gamma dimension bound and the
distinct Adams eigenvalues make this a class in
$F_\gamma^2K_0(E\times T)_{\mathbf Q}$, with the same component
restrictions. Its two fibres therefore give a relation in
$J_K(E)$ whose component image is $([x]-[y],0)$.

Every degree-zero class on $S$ is generated by such differences
with $x,y$ outside $C$. To justify moving the support, take a
smooth projective curve through any unwanted point and not
contained in $C$, then use Riemann--Roch on that curve to replace
its point divisor by a divisor supported in the complementary
open set. This uses cycle representatives, not a nonexistent
pushforward from the Chow group of an open immersion.

\emph{2. Kill the entire double-curve ambiguity by proper families.}
At a point parameter, the weight-two part of (FN.6) gives
\begin{equation}
 K_1(C)_{\mathbf Q}^{(2)}\xrightarrow{\delta}S(E)
 \longrightarrow\mathrm{CH}_0(S)_{\mathbf Q}\oplus
                    \mathrm{CH}_0(R)_{\mathbf Q}
 \longrightarrow0.                                  \tag{FN.7}
\end{equation}
The last zero follows from $F_\gamma^2K_0(C)=0$; rational
weight projection makes the preceding sequence exact in this
weight. We check that the first image is contained in $J_K(E)$.

The Gersten sequence on the smooth curve shows that
$K_1(C)^{(2)}$ is generated, rationally, by
$u[\mathcal O_c]$ with $u\in k^*$ and $c\in C(k)$.
Indeed the generic-point $K_1(k(C))$ has weight one, so a
weight-two class lies in the point-pushforward subgroup.
Write
\[
 u[\mathcal O_c]
   =u\bigl([\mathcal O_c]-[\mathcal O_{z_1}]\bigr)
                         +u[\mathcal O_{z_1}].
\]
The first summand is the difference of two fibres of an actual
Poincare line-bundle family on $C$ with a smooth proper curve
parameter. Multiply its rank-zero $K_0$ class by the constant
$K_1$ unit $u$, apply the natural Mayer--Vietoris boundary on
the whole product, and then apply the weight-two projector on
$K_0$ of the product. This produces exactly its boundary as a
relation in $J_K(E)$. One may use $C$ itself with the normalized
diagonal Poincare line bundle, or a smooth curve in its Jacobian.

The second summand is already in the restriction image from
$K_1(R)$. In fact the rank-zero line-bundle class
$[\mathcal O_R(e_1)]-1$ restricts to
$[\mathcal O_C(z_1)]-1=[\mathcal O_{z_1}]$, and multiplying
by the constant unit gives the asserted lift. Its boundary
in (FN.7) is zero. There are no triple strata in this two-component
argument, so no additional compatibility is being suppressed.

\emph{3. Take the quotient.}
The blow-up formula gives
$\mathrm{CH}_0(R)=\mathbf Z$. Step~1 supplies every degree-zero
relation on the first component, while Step~2 removes the full
kernel of the component map. Degrees themselves are constant
in all proper perfect families. Thus (FN.7) proves (FN.2).
For an actual exceptional realization satisfying EP, the
negative-definite degree matrix in C.3bm makes exceptional
pushforwards span both remaining coordinates. EP then gives
$B^2(X)_{\mathbf Q}=0$.

\emph{Application and limit.}
This applies intrinsically to the proposed plane--degree-$q$
surface test whenever its asserted regular resolution has
exactly $S$ and $R=\operatorname{Bl}_{C\cap G}\mathbf P^2$
as components: take $A=\mathcal O_S(1)$, $h=H$, and
$g=G|_S$, of degree $q+2$. In particular the quintic test has
$35$ such points. The argument settles the intrinsic family
obstruction for those data. It does not by itself certify the
proposed hypersurface resolution or its ampleness, and it is
not a proof for arbitrary SNC exceptional surfaces. Blowing
up a general double-curve point creates an additional plane
component and a triple point; the above two-component
Mayer--Vietoris argument cannot simply be reused without
checking that new global compatibility.
